\section{Введение}
\label{sec:Chapter0} \index{Chapter0}

За последние годы наблюдается быстрый прогресс в экспериментальной квантовой аппаратуре, в том числе в платформах на нейтральных атомах, управляемых оптическими пинцетами и возбуждаемых в ридберговские состояния \cite{saffman2010quantum,browaeys2020manybody}. Эти системы одновременно обладают масштабируемостью по числу кубитов и сильными многочастичными взаимодействиями Ван-дер-Ваальса, что делает их естественной основой для аналоговых (Hamiltonian-based) вычислений и квантового моделирования \cite{ebadi2021phases,wurtz2023aquila}.

Квантовое машинное обучение (QML) является одним из наиболее активно исследуемых направлений применения устройств промежуточного масштаба \cite{biamonte2017qml,schuld2021mlqc}. Однако для вариационных квантовых моделей, построенных из аппаратно-эффективных параметризованных схем (PQC/VQC), существуют как оптимизационные, так и выразительные ограничения: в частности, возможны явления плато градиентов (barren plateaus) \cite{mcclean2018barren,cerezo2021vqa}, а классическая аппроксимация таких моделей в ряде режимов оказывается эффективной. Существенный класс результатов связывает выход PQC с частичными рядами Фурье по входным данным, причём доступный спектр частот определяется спектром кодирующего гамильтониана \cite{schuld2020encoding}. На этой основе показано, что вариационные QML-модели могут быть приближены (а иногда и «деквантованы») классическими суррогатами на основе случайных фурье-признаков (random Fourier features) \cite{landman2022rff,sweke2025rfflimits}. Следовательно, при проектировании QML-анзаца важно опираться на физические динамики, индуцирующие богатый спектр и многочастичные корреляции, менее подверженные таким классическим приближениям \cite{schuld2020encoding,sweke2025rfflimits}.

Ридберговские массивы нейтральных атомов предоставляют именно такой тип динамики: эволюция задаётся многочастичным гамильтонианом с дальнодействующим взаимодействием и может программироваться через временные профили частоты Раби, фазы и (в ряде режимов) локальных детюнингов \cite{saffman2010quantum,ebadi2022mis,wurtz2023aquila}. В контексте QML на таких устройствах рассматриваются различные подходы. Один из них --- квантовое резервуарное вычисление, где квантовая динамика используется как фиксированный высокоразмерный «резервуар», а обучается лишь классический считывающий слой \cite{fujii2017qrc}. В данной работе мы исследуем альтернативную постановку: квантовую нейронную сеть (QNN), в которой параметры аналоговой эволюции (импульсы управления) являются обучаемыми.

Переход к обучаемым аналоговым QNN на ридберговском устройстве требует решения двух ключевых задач: (i) построения физически реализуемой параметризации управляющих импульсов, совместимой с ограничениями платформы; (ii) разработки численно эффективного метода симуляции промежуточных размеров системы с возможностью вычисления градиентов для обучения.

В работе предложен параметризованный QNN-анзац для ридберговской платформы и реализован ускоренный численный симулятор на базе JAX \cite{bradbury2018jax}, позволяющий выполнять батчевую симуляцию и дифференцирование по параметрам импульсов. Эффективность симуляции достигается за счёт аппроксимации временной зависимости кусочно-постоянными участками и использования редукции пространства состояний в режиме ридберговской блокады. На ряде синтетических задач классификации демонстрируется обучаемость модели и анализируется вклад межатомного взаимодействия и измеряемых корреляторов.

Структура работы следующая. В разделе 2 выводится многочастичный гамильтониан ридберговской системы и описывается параметризация QNN. В разделе 3 рассматриваются численные методы симуляции, аппроксимации и ускорения, включая анализ ридберговской блокады. В разделе 4 приведены результаты обучения, абляционные эксперименты и анализ гиперпараметров, подобранных с помощью Optuna \cite{akiba2019optuna}.
